\documentclass[12pt]{article}
\usepackage[utf8]{inputenc}
\usepackage[margin=0.75in]{geometry}
\usepackage{amsmath}
\usepackage{amssymb}
\usepackage{amsthm}
\usepackage{graphicx} % Required for inserting images
\usepackage{array}
\usepackage{tikz-cd}


\title{MAT1100 Homework 1}
\author{Aaron Avram}
\date{September 19 2026}

\begin{document}
\newtheorem{definition}{Definition}
\newtheorem{theorem}{Theorem}
\newtheorem{claim}{Claim}
\newtheorem{corollary}{Corollary}

\maketitle

\begin{align*}
    \textbf{Question 3}
\end{align*}

\begin{enumerate}
    \item[(a)]
    \begin{proof}
        First consider $x \in C_{D_8}(r)$. Since $D_8$ is generated by $r$ and $s$, there exist $a, b \in \mathbb Z$
        such that $x = r^as^b$. Additionally $xr = rx$ which implies:
        \[ r^as^br = rr^as^b \]
        simplifying we get
        \[ s^br = rs^b \]
        now without loss of generality $b \in \{ 0, 1\}$
        and if $b = 1$ this would imply
        \[ sr = rs \]
        which contradicts the group presentation
        hence $b = 0$. Therefore
        \[ C_{D_8}(r) = \langle r \rangle \]

        Next consider $x \in C_{D_8}(r)$. $x = r^as^b$ for $a, b \in \mathbb Z$.
        and since $xs = sx$ we get
        \[ r^as^bs = sr^as^b \]
        and using $sr = r^{-1}s$ we find
        \[ r^as^bs = r^{-a}s^bs \]
        using cancellation
        \[ r^a = r^{-a} \implies r^{2a} = 0 \]
        which implies $a = 2 \operatorname{mod} 4$, i.e. $a$ is even. 
        Therefore
        \[ C_{D_8}(s) = \langle s, r^2 \rangle \]
    \end{proof}

    \item[(b)]
    \begin{proof}
        First note that $Z(D_8) = C_{D_8}(r) \cap C_{D_8}(s)$. This follows as if an element is in the center
        it commutes with all of $D_8$ and in particular $r$ and $s$; conversely, since $r$ and $s$ generate $D_8$
        if an element of $D_8$ commutes with both of them it commutes with the whole group and is hence in the center.
        Therefore
        \[ Z(D_8) = \langle e, r^2 \rangle \]
        now to determine $Cl_{D_8}(r)$ first note that by orbit-stabilizer (applied to $D_8 \curvearrowright D_8$ by conjugation)
        \[ |Cl_{D_8}(r)| = \frac{|D_8|}{C_{D_8}(r)} = \frac{8}{4} = 2 \]
        then by the presentation we have $sr = r^{-1}s$ and since $r^{-1} = r^3$ we have
        \[ srs^{-1} = r^{-1} = r^3 \]
        and hence $Cl_{D_8}(r) = \{ r, r^3 \}$. Similarly
        \[ |Cl_{D_8}(s)| = 2 \]
        and by inspection
        \[ rsr^{-1} = r^2s \]
        therefore
        \[ Cl_{D_8}(s) = \{ s, r^2s \} \]

        Finally since the center and non-trivial conjugacy classes of $D_8$ partition it,
        this implies that the remaining conjugacy classes have size 2. However the remaining elements
        are not central, and this means their conjugacy classes have size at least 2. Thus there is
        only one conjugacy class remaining, and it is
        \[ Cl_{D_8}(rs) = \{rs, r^3s \} \]
    \end{proof}
    \item[(c)]
    \begin{proof}
        the class equation for $G$ is:
        \[ |G| = |Z(G)| + |Cl_{D_8}(r)| + |Cl_{D_8}(s)| + |Cl_{D_8}(rs)| \]
        with
        \[ Z(G) = \{ e, r^2 \} \]
        \[ Cl_{D_8}(r) = \{ r, r^3 \} \]
        \[ Cl_{D_8}(s) = \{ s, r^2s \} \]
        \[ Cl_{D_8}(rs) = \{ rs, r^3s \} \]
    \end{proof}
\end{enumerate}

\end{document}